% GENERATED FILE. Edit canonical lessons, then run npm run generate:books.
% canonical-source-hash: fnv1a64:2a694916ed88f2a6
% canonical-lessons: TA-C496-karppam, TA-C496-piracavam, TA-C496-pira, TA-C496-pillai, TA-C496-tattetu

\chapter{Birth and Adoption}
\label{ch:ta-birth-and-adoption}
\hlchaptermodality{496}

\begin{chapteropening}
\textbf{By the end of this chapter:} \emph{I can talk about pregnancy, childbirth, being born, children and adoption.}

Complete the last lesson of chapter 496: I can talk about pregnancy, childbirth, being born, children and adoption.
\end{chapteropening}

\section[\texorpdfstring{\ta{கர்ப்பம்} (karppam)}{கர்ப்பம் (karppam)}]{\ta{கர்ப்பம்} (karppam) — pregnancy}

\label{lesson:TA-C496-karppam}



\begin{quote}
Before the new one: say the Tamil for to sing a lullaby, then the Tamil for to feed.
\end{quote}

\subsection*{You'll want to know: \ta{கர்ப்பம்}}
\textbf{\ta{கர்ப்பம்}} — \emph{karppam} — \textquotedblleft{}pregnancy\textquotedblright{}.

\begin{quote}
\ta{என்} \ta{அக்கா} \ta{இப்போது} \ta{ஆறு} \ta{மாதக்} \ta{கர்ப்பமாக} \ta{இருக்கிறாள்}.
\end{quote}

\begin{quote}
\emph{eṉ akkā ippōtu āṟu mātak karppamāka irukkiṟāḷ.}
\end{quote}

\begin{quote}
My elder sister is six months pregnant now.
\end{quote}

\ta{கர்ப்பமாக} \ta{இரு}, to be pregnant.

\begin{cousinweb}
\textbf{\ta{கர்ப்பம்}} (\emph{karppam}) is borrowed from Sanskrit \emph{garbha}, \textquotedblleft{}womb, foetus\textquotedblright{} (CDIAL 4055); Tamil writes the \emph{g} and \emph{bh} with \ta{க} and \ta{ப}.
\end{cousinweb}

\begin{grammarlens}[title={What you've built}]
One more word for describing what you lived and felt.
\end{grammarlens}

\subsection*{Your turn}
\begin{itemize}
\raggedright
  \item \emph{Say it:} \emph{karppam}
  \item \emph{Say it:} \emph{karppam}, once more
  \item \emph{Say it:} \emph{eṉ akkā ippōtu āṟu mātak karppamāka irukkiṟāḷ.} — My elder sister is six months pregnant now.
  \item \emph{Recall:} say the Tamil for to sing a lullaby, then the Tamil for to feed, then say \emph{karppam} again
\end{itemize}

\begin{tcolorbox}[breakable,colback=teal!4,colframe=teal!35!black,title={Before you move on}]
What does \ta{கர்ப்பம்} mean? (\textquotedblleft{}Pregnancy\textquotedblright{}.) Say it once more.
\end{tcolorbox}
\section[\texorpdfstring{\ta{பிரசவம்} (piracavam)}{பிரசவம் (piracavam)}]{\ta{பிரசவம்} (piracavam) — childbirth; delivery}

\label{lesson:TA-C496-piracavam}



\begin{quote}
Before the new one: say the Tamil for to feed, then the Tamil for pregnancy.
\end{quote}

\subsection*{You'll want to know: \ta{பிரசவம்}}
\textbf{\ta{பிரசவம்}} — \emph{piracavam} — \textquotedblleft{}childbirth; delivery\textquotedblright{}.

\begin{quote}
\ta{அவளுக்கு} \ta{அரசு} \ta{மருத்துவமனையில்} \ta{பிரசவம்} \ta{நடந்தது}.
\end{quote}

\begin{quote}
\emph{avaḷukku aracu maruttuvamaṉaiyil piracavam naṭantatu.}
\end{quote}

\begin{quote}
She gave birth in the government hospital.
\end{quote}

\ta{பிரசவ} \ta{வலி}, labour pains.

\begin{cousinweb}
\textbf{\ta{பிரசவம்}} (\emph{piracavam}) is borrowed from Sanskrit \emph{prasava}, \textquotedblleft{}procreation, birth\textquotedblright{} (CDIAL 8827).
\end{cousinweb}

\begin{grammarlens}[title={What you've built}]
One more word for describing what you lived and felt.
\end{grammarlens}

\subsection*{Your turn}
\begin{itemize}
\raggedright
  \item \emph{Say it:} \emph{piracavam}
  \item \emph{Say it:} \emph{piracavam}, once more
  \item \emph{Say it:} \emph{avaḷukku aracu maruttuvamaṉaiyil piracavam naṭantatu.} — She gave birth in the government hospital.
  \item \emph{Recall:} say the Tamil for to feed, then the Tamil for pregnancy, then say \emph{piracavam} again
\end{itemize}

\begin{tcolorbox}[breakable,colback=teal!4,colframe=teal!35!black,title={Before you move on}]
What does \ta{பிரசவம்} mean? (\textquotedblleft{}Childbirth; delivery\textquotedblright{}.) Say it once more.
\end{tcolorbox}
\section[\texorpdfstring{\ta{பிற} (piṟa)}{பிற (piṟa)}]{\ta{பிற} (piṟa) — to be born}

\label{lesson:TA-C496-pira}



\begin{quote}
Before the new one: say the Tamil for pregnancy, then the Tamil for childbirth.
\end{quote}

\subsection*{You'll want to know: \ta{பிற}}
\textbf{\ta{பிற}} — \emph{piṟa} — \textquotedblleft{}to be born\textquotedblright{}.

\begin{quote}
\ta{நான்} \ta{ஒரு} \ta{சிறிய} \ta{கிராமத்தில்} \ta{பிறந்தேன்}.
\end{quote}

\begin{quote}
\emph{nāṉ oru ciṟiya kirāmattil piṟantēṉ.}
\end{quote}

\begin{quote}
I was born in a small village.
\end{quote}

Past (I …): \ta{பிறந்தேன்}, \emph{piṟantēṉ}.

Past with -\ta{ந்த}-: \ta{பிறந்தேன்}. \ta{பிறந்தநாள்}, a birthday, is the born-day.

\begin{cousinweb}
\textbf{\ta{பிற}} (\emph{piṟa}) is native Dravidian (DEDR 4422), in the entry of \ta{பெறு} \emph{peṟu}, \textquotedblleft{}to get, to bear\textquotedblright{}.
\end{cousinweb}

\begin{grammarlens}[title={What you've built}]
One more word for describing what you lived and felt.
\end{grammarlens}

\subsection*{Your turn}
\begin{itemize}
\raggedright
  \item \emph{Say it:} \emph{piṟa}
  \item \emph{Say it:} \emph{piṟa}, once more
  \item \emph{Say it:} \emph{nāṉ oru ciṟiya kirāmattil piṟantēṉ.} — I was born in a small village.
  \item \emph{Recall:} say the Tamil for pregnancy, then the Tamil for childbirth, then say \emph{piṟa} again
\end{itemize}

\begin{tcolorbox}[breakable,colback=teal!4,colframe=teal!35!black,title={Before you move on}]
What does \ta{பிற} mean? (\textquotedblleft{}To be born\textquotedblright{}.) Say it once more.
\end{tcolorbox}
\section[\texorpdfstring{\ta{பிள்ளை} (piḷḷai)}{பிள்ளை (piḷḷai)}]{\ta{பிள்ளை} (piḷḷai) — a son or daughter; offspring}

\label{lesson:TA-C496-pillai}



\begin{quote}
Before the new one: say the Tamil for childbirth, then the Tamil for to be born.
\end{quote}

\subsection*{You'll want to know: \ta{பிள்ளை}}
\textbf{\ta{பிள்ளை}} — \emph{piḷḷai} — \textquotedblleft{}a son or daughter; offspring\textquotedblright{}.

\begin{quote}
\ta{அவர்களுக்கு} \ta{மூன்று} \ta{பிள்ளைகள்}.
\end{quote}

\begin{quote}
\emph{avarkaḷukku mūṉṟu piḷḷaikaḷ.}
\end{quote}

\begin{quote}
They have three children.
\end{quote}

Plural: \ta{பிள்ளைகள்}, \emph{piḷḷaikaḷ}.

Children of a family, sons or daughters: \ta{பிள்ளைகள்}.

\begin{cousinweb}
\textbf{\ta{பிள்ளை}} (\emph{piḷḷai}) is native Dravidian (DEDR 4198), for a child and the young of many animals.
\end{cousinweb}

\begin{grammarlens}[title={What you've built}]
One more word for describing what you lived and felt.
\end{grammarlens}

\subsection*{Your turn}
\begin{itemize}
\raggedright
  \item \emph{Say it:} \emph{piḷḷai}
  \item \emph{Say it:} \emph{piḷḷai}, once more
  \item \emph{Say it:} \emph{avarkaḷukku mūṉṟu piḷḷaikaḷ.} — They have three children.
  \item \emph{Recall:} say the Tamil for childbirth, then the Tamil for to be born, then say \emph{piḷḷai} again
\end{itemize}

\begin{tcolorbox}[breakable,colback=teal!4,colframe=teal!35!black,title={Before you move on}]
What does \ta{பிள்ளை} mean? (\textquotedblleft{}A son or daughter; offspring\textquotedblright{}.) Say it once more.
\end{tcolorbox}
\section[\texorpdfstring{\ta{தத்தெடு} (tatteṭu)}{தத்தெடு (tatteṭu)}]{\ta{தத்தெடு} (tatteṭu) — to adopt}

\label{lesson:TA-C496-tattetu}



\begin{quote}
Before the new one: say the Tamil for to be born, then the Tamil for a son or daughter.
\end{quote}

\subsection*{You'll want to know: \ta{தத்தெடு}}
\textbf{\ta{தத்தெடு}} — \emph{tatteṭu} — \textquotedblleft{}to adopt\textquotedblright{}.

\begin{quote}
\ta{அவர்கள்} \ta{ஒரு} \ta{பெண்} \ta{குழந்தையைத்} \ta{தத்தெடுத்தார்கள்}.
\end{quote}

\begin{quote}
\emph{avarkaḷ oru peṇ kuḻantaiyait tatteṭuttārkaḷ.}
\end{quote}

\begin{quote}
They adopted a baby girl.
\end{quote}

Past (I …): \ta{தத்தெடுத்தேன்}, \emph{tatteṭuttēṉ}.

\ta{தத்து} \textquotedblleft{}adoption\textquotedblright{} + \ta{எடு} \textquotedblleft{}to take\textquotedblright{}.

\begin{cousinweb}
\textbf{\ta{தத்தெடு}} (\emph{tatteṭu}) is a compound of \ta{தத்து}, \textquotedblleft{}adoption\textquotedblright{}, borrowed from Sanskrit \emph{datta}, \textquotedblleft{}given\textquotedblright{} (CDIAL 6140), and native \ta{எடு}, \textquotedblleft{}to take\textquotedblright{}.
\end{cousinweb}

\begin{grammarlens}[title={What you've built}]
One more word for describing what you lived and felt.
\end{grammarlens}

\subsection*{Your turn}
\begin{itemize}
\raggedright
  \item \emph{Say it:} \emph{tatteṭu}
  \item \emph{Say it:} \emph{tatteṭu}, once more
  \item \emph{Say it:} \emph{avarkaḷ oru peṇ kuḻantaiyait tatteṭuttārkaḷ.} — They adopted a baby girl.
  \item \emph{Recall:} say the Tamil for to be born, then the Tamil for a son or daughter, then say \emph{tatteṭu} again
\end{itemize}

\begin{tcolorbox}[breakable,colback=teal!4,colframe=teal!35!black,title={Before you move on}]
What does \ta{தத்தெடு} mean? (\textquotedblleft{}To adopt\textquotedblright{}.) Say it once more.
\end{tcolorbox}
